\section{The extreme roots of Cohen's Lambert polynomials}
\label{sec:cohen-edges}

The diagonal polynomials are a previously studied family.  In the
notation of Cohen~\cite[Proposition~4.8]{Cohen2021},
\begin{equation}\label{eq:cohen-identification}
 L_n(X)=(-1)^n p_n(X).
\end{equation}
Cohen's Conjecture~7.1 concerns their real roots, interlacing, and
extreme-root asymptotics, including a numerical constant in the largest
root.  We prove a full fixed-index extreme-root expansion, establish
the conjectured polynomial dependence of every coefficient on the
integer parameter, and identify the first constant.  The proof
also applies to Cohen's generalized polynomials
$L_{n,k}$ for each fixed positive integer $k$.  All references to that
conjecture below concern the version dated 18 January 2021; no claim is
made that every aspect of the conjecture remains unresolved elsewhere.

\subsection{An operator normalization and a reciprocal transform}

Let $n,k\geq1$, set $N=n+k-1$, and put $D=d/dx$.  Cohen's signed-Stirling
formula is equivalently
\begin{equation}\label{eq:cohen-operator}
 L_{n,k}(X)=\frac{(-1)^n k}{n}
   \left.\left(\prod_{r=1}^{n-1}(1+D/r)\right)x^N\right|_{x=-X}.
\end{equation}
An empty product is one.  For $k=1$, this gives
\eqref{eq:cohen-identification}.  The operator argument in
Lemma~\ref{lem:zero-law-inputs} proves immediately that $L_{n,k}$ has
a zero of multiplicity $k$ at the origin and exactly $n-1$ additional
zeros, all simple and positive.  Indeed each application of $1+cD$,
$c>0$, removes one multiplicity at zero and introduces one negative
simple root, with the other negative roots strictly interlacing those
of the preceding polynomial.  Thus this elementary argument proves
the real-rooted portions of Cohen's Conjecture~7.1(1),(4).

\begin{proposition}[Interlacing of successive Lambert polynomials]
\label{prop:cohen-interlacing}
The zeros of $L_n$ and $L_{n+1}$ interlace.  Their only common zero
is $0$, and all interlacing inequalities between positive zeros are
strict.
\end{proposition}

\begin{proof}
Write $T_n=\prod_{r=1}^{n-1}(1+D/r)$ and $P_n(x)=T_nx^n$.
For every real $t$,
\[
 P_{n+1}(x)+tP_n(x)
    =T_n\bigl[x^n(x+(n+1)/n+t)\bigr]
\]
is real-rooted, because its input polynomial is real-rooted and
each factor $1+D/r$ preserves that property.  The elementary
real-rooted pencil criterion implies weak interlacing.
To recall why that criterion applies, cancel common factors in
$R/P$, where every $R+tP$ is real-rooted.  This real rational
function cannot take a real value in the upper half-plane.  Its
imaginary part therefore has constant sign there; the monic degree
difference one makes that sign positive.  Its residues at the
remaining real poles are negative, which forces strict interlacing
of the remaining zeros and poles.  Reintroducing common factors
gives weak interlacing.

In the present case no nonzero root can be shared.  The identity
\[
 P_{n+1}'=(n+1)(P_n+P_n'/n)
\]
would make a shared nonzero root a multiple root of
$P_{n+1}-(n+1)P_n/n=T_nx^{n+1}$.  The operator argument proves that
this last polynomial has only simple nonzero roots, with zero
of multiplicity two, a contradiction.  Changing $x$ to $-X$
proves the assertion.  This also proves the interlacing assertion
in Cohen's Conjecture~7.1(3).
\end{proof}

For $1\leq j\leq n-1$, write $X_{n,k,j}$ for the positive roots in
\emph{decreasing} order.  This indexing is intentionally opposite to
the bulk indexing in Section~\ref{sec:velocity-zero-law}.
Let $e_r=e_r(1,1/2,\ldots,1/(n-1))$ be elementary symmetric functions,
with $e_0=1$, and use $N^{\underline r}=N(N-1)\cdots(N-r+1)$.
The reciprocal polynomial
\begin{equation}\label{eq:reciprocal-lambert-polynomial}
 Q_{n,N}(w)=\sum_{r=0}^{n-1}(-1)^r
       \frac{N^{\underline r}}{N^r}e_r w^r
\end{equation}
has its $n-1$ roots at $w=N/X_{n,k,j}$, in increasing order.
It is obtained from $L_{n,k}(N/w)$ by multiplication by
$n(w/N)^N/((-1)^{k-1}k)$, so the transform introduces no finite
nonzero roots.  Define also
\begin{equation}\label{eq:gamma-product-fn}
 f_n(w)=\prod_{r=1}^{n-1}(1-w/r)
       =\frac{\Gamma(n-w)}{\Gamma(n)\Gamma(1-w)},\qquad
 g(w)=\frac1{\Gamma(1-w)}.
\end{equation}
At the apparent poles the gamma quotient is understood through the
polynomial product.  Coefficient extraction gives the exact identity
\begin{equation}\label{eq:reciprocal-coefficient-integral}
 Q_{n,N}(w)=\frac{N!}{N^N}[z^N]\,e^{Nz}f_n(wz).
\end{equation}

\subsection{An exact logarithmic shift and a uniform saddle expansion}

The important simplification is to remove $\log n$ \emph{before}
expanding a zero.  Set $\ell=\log n$ and define
\begin{equation}\label{eq:log-shifted-transform}
 h_n(v)=n^v f_n(v),\qquad
 T_{n,N}(v)=\frac{N!}{N^N}[z^N]e^{Nz}h_n(vz).
\end{equation}
Combining exponentials in \eqref{eq:reciprocal-coefficient-integral}
and rescaling its coefficient variable gives the exact identity
\begin{equation}\label{eq:exact-logarithmic-shift}
 Q_{n,N}(w)=r^N T_{n,N}(v),\qquad
 r=1-\frac{w\ell}{N},\qquad v=\frac wr.
\end{equation}
Indeed $e^{Nz}f_n(wz)=e^{(N-w\ell)z}h_n(wz)$, and replacing
$z$ by $z/r$ proves the formula.  At a root corresponding to
$X=N/w$, the new coordinate is simply
\begin{equation}\label{eq:shifted-reciprocal-coordinate}
 v=\frac{N}{X-\log n},\qquad X=\log n+\frac Nv.
\end{equation}
All these transformations are used near a fixed positive $v$, where
$r$ is nonzero for large $n$.  Zeros far from this region need not
be classified to determine any fixed number of largest roots.

The uniform gamma-ratio expansion \cite[\S5.11]{DLMFGamma} gives,
locally uniformly in the complex plane and with every fixed number
of derivatives,
\begin{equation}\label{eq:gamma-ratio-full}
 h_n(v)\sim g(v)R(v;1/n),\qquad
 R(v;\epsilon)=\exp\left(\sum_{a\geq1}\lambda_a(v)\epsilon^a\right),
\end{equation}
where
\begin{equation}\label{eq:gamma-ratio-Bernoulli}
 \lambda_a(v)=\frac{(-1)^{a+1}}{a(a+1)}
      \bigl(B_{a+1}(-v)-B_{a+1}(0)\bigr).
\end{equation}
Here $B_m(x)$ are the Bernoulli polynomials defined by
$t e^{xt}/(e^t-1)=\sum_{m\geq0}B_m(x)t^m/m!$.
The symbol $\sim$ denotes an asymptotic expansion to every finite
order; no convergence of the formal series $R(v;\epsilon)$ is
assumed.  Its first coefficient is
\begin{equation}\label{eq:gamma-ratio-edge}
 R(v;\epsilon)=1+q_1(v)\epsilon+O(\epsilon^2),\qquad
 q_1(v)=\frac{v(v+1)}2.
\end{equation}
In particular $h_n$ and each of its fixed derivatives are uniformly
bounded on every fixed compact set.  The gamma quotient is analytic
on a fixed compact set for large $n$, and its expansion there follows
from uniform Stirling expansion; derivative estimates follow by
Cauchy's formula on a slightly larger compact set.

\begin{lemma}[An all-orders coefficient-extraction estimate]
\label{lem:reciprocal-saddle}
For $r\geq0$, define polynomials $m_r(\eta)$ by
\begin{equation}\label{eq:saddle-moment-generating-function}
 \sum_{r\geq0}m_r(\eta)\frac{s^r}{r!}
  =\exp\left(\sum_{q\geq2}
        \frac{(-1)^{q-1}}q s^q\eta^{q-1}\right),
\end{equation}
and differential operators
\begin{equation}\label{eq:saddle-operator-definition}
 \mathcal D_b F(v)=\sum_{r=0}^{2b}
       [\eta^b]m_r(\eta)\,\frac{v^r}{r!}F^{(r)}(v).
\end{equation}
Then, for each fixed integer $M\geq0$ and compact $K\subset\mathbb C$,
\begin{equation}\label{eq:log-free-saddle-expansion}
 T_{n,N}(v)=\sum_{b=0}^{M}N^{-b}\mathcal D_bh_n(v)
             +O_{K,M}(N^{-M-1})
\end{equation}
uniformly for $v\in K$ and all sufficiently large independent
positive integers $n,N$.  In particular,
\begin{align}\label{eq:first-saddle-operators}
 \mathcal D_0F&=F,\qquad
 \mathcal D_1F=-\frac{v^2}{2}F'',\\
 \mathcal D_2F&=\frac{v^3}{3}F'''+\frac{v^4}{8}F^{(4)},\notag\\
 \mathcal D_3F&=-\frac{v^4}{4}F^{(4)}
                 -\frac{v^5}{6}F^{(5)}
                 -\frac{v^6}{48}F^{(6)}.\notag
\end{align}
Consequently $T_{n,N}\to g$ locally uniformly as $n,N\to\infty$.
\end{lemma}

\begin{proof}
On $z=e^{i\varphi}$ the exact coefficient integral is
\[
 \frac{N!}{2\pi N^N}\int_{-\pi}^{\pi}
       e^{N e^{i\varphi}-iN\varphi}h_n(ve^{i\varphi})\,d\varphi.
\]
After extracting $e^N$, Stirling's formula bounds its prefactor by
$O(\sqrt N)$, while the modulus of its exponential factor is
$e^{N(\cos\varphi-1)}$.  On $|\varphi|\leq N^{-2/5}$, expand the
amplitude about $z=1$ through degree $2M+1$.  The uniform compact
bounds on $h_n^{(2M+2)}$ bound the Taylor remainder by
$C_{K,M}|e^{i\varphi}-1|^{2M+2}$, independently of large $n$.
Its integral is $O_{K,M}(N^{-M-1})$, since
\[
 \sqrt N\int e^{-cN\varphi^2}|\varphi|^{2M+2}\,d\varphi
            =O_M(N^{-M-1}).
\]
On the complementary arc the exponential is at most
$e^{-cN^{1/5}}$.  The amplitude and its finite Taylor polynomial
remain uniformly bounded on a fixed compact set, so this arc is
negligible.  We may therefore integrate the Taylor polynomial over
the whole circle.

Its exact moments are
\[
 M_r(N)=\frac{N!}{N^N}[z^N]e^{Nz}(z-1)^r
       =\sum_{a=0}^r(-1)^{r-a}\binom ra
            \frac{N^{\underline a}}{N^a}.
\]
Their exponential generating function is exactly
\[
 \sum_{r\geq0}M_r(N)\frac{s^r}{r!}
       =e^{-s}(1+s/N)^N
       =\exp\left(\sum_{q\geq2}
               \frac{(-1)^{q-1}s^q}{qN^{q-1}}\right).
\]
Thus $M_r(N)=m_r(1/N)$.  Moreover $m_0=1$, $m_1=0$, and for
$r\geq2$ the least possible power of $\eta$ is $\lceil r/2\rceil$:
a product of terms $s^q\eta^{q-1}$ has $s$-degree at most twice
its $\eta$-degree.  This proves both the finite range in
\eqref{eq:saddle-operator-definition} and the fact that the
Taylor term of degree $2M+1$ contributes only $O(N^{-M-1})$.
Grouping the other terms by powers of $1/N$ proves
\eqref{eq:log-free-saddle-expansion}.  The displayed operators
follow directly from the generating function.
\end{proof}

\subsection{All-orders edge asymptotics and polynomial parameter dependence}

Write $H_m^{(r)}=\sum_{a=1}^m a^{-r}$, $H_0^{(r)}=0$, and
$H_m=H_m^{(1)}$.  Let $\gamma$ be Euler's constant.

\begin{theorem}[All-orders fixed-index largest roots]
\label{thm:cohen-extreme-roots}
Fix positive integers $j,k$.  There are explicitly computable
polynomials $C_{j,m}(\kappa)$, $m\geq1$, with
$\deg C_{j,m}=m-1$, such that for every fixed integer $M\geq0$,
\begin{equation}\label{eq:cohen-extreme-root-all-orders}
 X_{n,k,j}=\frac{n+k-1}{j}+\log n+\gamma-H_{j-1}
     +\sum_{m=1}^{M}\frac{C_{j,m}(k-1)}{n^m}
     +O_{j,k,M}(n^{-M-1}).
\end{equation}
The sum is empty when $M=0$.  Put
\[
 s_2=H_{j-1}^{(2)}+\zeta(2),\qquad
 s_3=H_{j-1}^{(3)}-\zeta(3).
\]
The first coefficient, independent of $\kappa$, is
\begin{equation}\label{eq:cohen-Cj}
 C_{j,1}(\kappa)=j^2s_3-js_2-j-\frac12=:C_j.
\end{equation}
For $m\geq2$, the coefficient of $\kappa^{m-1}$ in
$C_{j,m}(\kappa)$ is
\begin{equation}\label{eq:cohen-leading-parameter-coefficient}
 (-1)^{m-1}(j^2s_3-js_2),
\end{equation}
which is nonzero.  In particular, for every fixed $k\geq1$,
\begin{equation}\label{eq:cohen-largest-root}
 X_{n,k,1}=n+k-1+\log n+\gamma
     -\frac{\frac32+\zeta(2)+\zeta(3)}{n}+O_k(n^{-2}).
\end{equation}
\end{theorem}

\begin{proof}
First identify the relevant roots.  The limit in
Lemma~\ref{lem:reciprocal-saddle} is
$g(v)=1/\Gamma(1-v)$, whose zeros are the positive integers,
all simple.  Rouch\'e's theorem on small disjoint circles about
$1,\ldots,j$ and uniform nonvanishing on the intervening compact
real intervals, including the interval starting at $0$, show that
$T_{n,N}$ has one root near each of $1,\ldots,j$ and no additional
smaller positive roots.  The nearby roots are real, either by
conjugation and uniqueness in each circle or by the real-rootedness
already established for the original polynomials.  For large $n$,
\eqref{eq:exact-logarithmic-shift} maps each of these roots to a
positive root $X=\log n+N/v$.  Positive $v$ correspond precisely to
roots with $X>\log n$ and preserve decreasing order in $X$.
Roots with $v<0$ give $X<\log n$ and cannot intervene among these
largest roots.  Hence the root approaching $j$ is exactly
\[
 v_{n,k,j}=\frac{N}{X_{n,k,j}-\log n}.
\]

We now use formal series only as a finite-order computational device.
Introduce independent formal variables $\epsilon,\eta$ and define
\begin{equation}\label{eq:formal-edge-equation}
 \mathcal T(v;\epsilon,\eta)
   =\sum_{b\geq0}\eta^b
           \mathcal D_b\bigl(g(v)R(v;\epsilon)\bigr).
\end{equation}
The gamma expansion and Lemma~\ref{lem:reciprocal-saddle} show that
this is the full two-parameter asymptotic expansion of $T_{n,N}$
when $\epsilon=1/n$ and $\eta=1/N$.  More precisely, its truncation
through total degree $d$ has error $O_{K,d}((1/n+1/N)^{d+1})$
on compact sets.  This follows by expanding $h_n$ to order $d-b$
in the term with $N^{-b}$ in
\eqref{eq:log-free-saddle-expansion}; the finitely many remainders
have exactly the asserted total order.

Since $g'(j)\ne0$, the formal implicit function theorem gives a
unique series $v_j(\epsilon,\eta)$ with constant term $j$ satisfying
$\mathcal T(v_j;\epsilon,\eta)=0$.  Crucially,
\[
 \mathcal T(v;\epsilon,0)=g(v)R(v;\epsilon).
\]
The factor $R$ is a formal unit and $g$ has its simple zero at $j$.
Thus $v_j(\epsilon,0)=j$, and $v_j-j$ is divisible by $\eta$.
The formal series
\begin{equation}\label{eq:formal-height-series}
 B_j(\epsilon,\eta)=\frac1\eta
          \left(\frac1{v_j(\epsilon,\eta)}-\frac1j\right)
\end{equation}
therefore contains only nonnegative powers of $\epsilon,\eta$.
All its coefficients are obtained recursively from
\eqref{eq:formal-edge-equation}.  Neither that series nor the
implicit series is asserted to converge.

These formal coefficients give genuine asymptotic expansions.
Indeed, restrict $N/n$ to a fixed compact subinterval of $(0,\infty)$,
which includes $N=n+k-1$ for fixed $k$ and large $n$.  Truncate the
formal root to any prescribed total degree and insert it into the
finite expansion of $T_{n,N}$.  The residual has the next total
order.  Locally uniform convergence $T_{n,N}\to g$, also for its
first derivative, and $g'(j)\ne0$ convert that residual into a root
error of the same order.  For real parameters this follows directly
from the mean value theorem.  Computing the root through total
order $M+1$, with residual $O(n^{-M-2})$, gives the height
$\log n+N/v$ with error $O(n^{-M-1})$.  This accounts explicitly
for the factor $N$ introduced by taking the reciprocal.

It remains to determine the polynomial degrees and the first
coefficients.  Set
\[
 c=H_{j-1}-\gamma.
\]
The local reciprocal-gamma expansion is
\begin{equation}\label{eq:reciprocal-gamma-at-integer}
 g(j+\delta)=g'(j)\delta
     \exp\left(c\delta-\frac{s_2}2\delta^2
                   +\frac{s_3}3\delta^3+O_j(\delta^4)\right),
 \qquad g'(j)=(-1)^j(j-1)!.
\end{equation}
It follows from
$g(j+\delta)=(-1)^j\delta\prod_{a=1}^{j-1}(a+\delta)/\Gamma(1-\delta)$
and the convergent Taylor expansion of $\log\Gamma(1-\delta)$.
In particular,
\[
 \frac{g''(j)}{g'(j)}=2c,\qquad
 \frac{g'''(j)}{g'(j)}=3(c^2-s_2),\qquad
 \frac{g^{(4)}(j)}{g'(j)}=4(c^3-3cs_2+2s_3).
\]
Use \eqref{eq:first-saddle-operators} and
\eqref{eq:gamma-ratio-edge} to solve
\eqref{eq:formal-edge-equation} through total degree two.  The result is
\begin{align}\label{eq:formal-reciprocal-root-two-terms}
 v_j(\epsilon,\eta)=j
   &+j^2c\eta+j^2q_1'(j)\epsilon\eta\\
   &+\bigl(j^3c^2+j^3s_2-j^4s_3\bigr)\eta^2
       +\text{terms of total degree at least three}.\notag
\end{align}
For example, the $\eta$ coefficient is
$-\mathcal D_1g(j)/g'(j)=j^2c$.  At $\epsilon=0$, the coefficient
of $\eta^2$ is
\[
 -\frac{\frac12g''(j)(j^2c)^2
      +(\mathcal D_1g)'(j)j^2c+\mathcal D_2g(j)}{g'(j)},
\]
which simplifies to the displayed value.  At first order in $\eta$,
replacing $g$ by $gR$ replaces $c$ by
$c+\partial_v\log R(j;\epsilon)$, giving the $\epsilon\eta$ term.
Consequently
\begin{equation}\label{eq:formal-height-first-coefficients}
 B_j(\epsilon,\eta)=-c
       -(j+\tfrac12)\epsilon+(j^2s_3-js_2)\eta
       +\text{terms of total degree at least two}.
\end{equation}

Now put $\kappa=k-1$ and substitute
\begin{equation}\label{eq:two-small-parameter-substitution}
 \eta=\frac{\epsilon}{1+\kappa\epsilon}.
\end{equation}
This defines the polynomials $C_{j,m}$ by
\begin{equation}\label{eq:Cjm-coefficient-recipe}
 B_j\left(\epsilon,\frac{\epsilon}{1+\kappa\epsilon}\right)
       =-c+\sum_{m\geq1}C_{j,m}(\kappa)\epsilon^m.
\end{equation}
Every nonconstant monomial $\epsilon^a\eta^b$ has $a+b\geq1$.
If $b\geq1$, its contribution to the coefficient of $\epsilon^m$
is a constant times $\kappa^{m-a-b}$, and if $b=0$ it is
independent of $\kappa$.  Therefore $\deg C_{j,m}\leq m-1$.
For $m\geq2$, the only possible term of degree $m-1$ comes from
$\eta$ itself; by \eqref{eq:formal-height-first-coefficients} its
coefficient is \eqref{eq:cohen-leading-parameter-coefficient}.
Since $s_3=-\sum_{a=j}^{\infty}a^{-3}<0$ and $s_2>0$, this
coefficient is nonzero.  At $m=1$ the constant polynomial is
$C_j=j^2s_3-js_2-j-1/2<0$, also nonzero.  This proves the exact
degree assertions and the stated formulas.  The formal-to-asymptotic
argument above completes \eqref{eq:cohen-extreme-root-all-orders}.
\end{proof}

\begin{corollary}[The first two largest-root corrections]
\label{cor:cohen-two-constants}
For every fixed positive integer $k$,
\begin{equation}\label{eq:cohen-two-corrections}
 X_{n,k,1}=n+k-1+\log n+\gamma
     +\frac{C_0(k)}n+\frac{C_1(k)}{n^2}+O_k(n^{-3}),
\end{equation}
where, in Cohen's indexing,
\begin{align}\label{eq:cohen-exact-constants}
 C_0(k)&=-\frac32-\zeta(2)-\zeta(3),\\
 C_1(k)&=(k-1)\bigl(\zeta(2)+\zeta(3)\bigr)
       +\zeta(2)^2-\zeta(2)\zeta(3)+2\zeta(3)+3\zeta(5)-\frac1{12}.
       \notag
\end{align}
Thus $C_0(k)=-4.346990970007820721872\ldots$ and
$C_1(1)=6.160067472396514377053\ldots$.
\end{corollary}

\begin{proof}
The first identity is \eqref{eq:cohen-Cj} at $j=1$.  For the second,
the gamma-ratio coefficient following $q_1$ is
$q_2(v)=v(v+1)(v+2)(3v+1)/24$.  Write, near $v=1$,
\[
 \frac{g(1+z)}{g'(1)}=z\exp\left(cz-\frac{s_2}2z^2
       +\frac{s_3}3z^3-\frac{s_4}4z^4+\frac{s_5}5z^5+O(z^6)\right),
\]
where $c=-\gamma$, $s_2=\zeta(2)$, $s_3=-\zeta(3)$,
$s_4=\zeta(4)$ and $s_5=-\zeta(5)$.  Substitute $q_1,q_2$ and
$\mathcal D_1,\mathcal D_2,\mathcal D_3$ into
\eqref{eq:formal-edge-equation}, then use
\eqref{eq:two-small-parameter-substitution}.  Equating three successive
coefficients in the simple-root equation and taking the reciprocal
in \eqref{eq:formal-height-series} gives
\[
 C_1(k)=-s_2^2+s_2s_3+(k-1)(s_2-s_3)-2s_3+5s_4-3s_5-\frac1{12}.
\]
The formal symbol $c$ cancels.  Since $5\zeta(4)=2\zeta(2)^2$,
this is \eqref{eq:cohen-exact-constants}.  This is a finite
coefficient calculation by the recursion already specified in
\eqref{eq:formal-edge-equation}; a rational-arithmetic replay is
included among the accompanying verification artifacts.  The
remainder follows from Theorem~\ref{thm:cohen-extreme-roots} with $M=2$.
\end{proof}

\begin{corollary}[Arithmetic form of every correction coefficient]
\label{cor:cohen-coefficient-algebra}
For fixed positive integers $j,m$,
\[
 C_{j,m}(\kappa)\in
   \mathbb Q[\zeta(2),\zeta(3),\ldots,\zeta(2m+1)][\kappa].
\]
Thus Euler's constant is confined to the leading shift in
\eqref{eq:cohen-extreme-root-all-orders}; it is absent from every
subsequent coefficient produced by the construction.
\end{corollary}

\begin{proof}
Fix $j$ and set $c=H_{j-1}-\gamma$.  Extract one more exponential,
writing $h_n(v)=e^{cv}a_n(v)$.  Exactly as in
\eqref{eq:exact-logarithmic-shift},
\[
 T_{n,N}(v)=(1+cv/N)^N
   \frac{N!}{N^N}[z^N]e^{Nz}a_n(sz),\qquad
 s=\frac{v}{1+cv/N}.
\]
The height is now $X=\log n-c+N/s$.  Up to the irrelevant nonzero
constant $e^{-cj}g'(j)$, the local asymptotic amplitude at $j$ is
\[
 \delta\exp\left(\sum_{r\geq2}
    \frac{(-1)^{r-1}H_{j-1}^{(r)}-\zeta(r)}r\delta^r\right)
       R(j+\delta;\epsilon).
\]
The finite harmonic sums are rational and the coefficients of $R$
are rational polynomials.  The recursion for a height coefficient
of order $n^{-m}$ uses the root through total degree $m+1$, hence
amplitude derivatives of order at most $2m+2$.  The initial zero
factor $\delta$ reduces the largest needed zeta index to $2m+1$.
After the displayed normalization, all divisions in the simple-root
recursion are by nonzero rational numbers.  This proves the claim.
\end{proof}

The theorem proves the fixed-$j$ assertion in Cohen's
Conjecture~7.1(2), since $H_n=\log n+\gamma+O(1/n)$, and proves
his full all-orders assertion in Conjecture~7.1(4) for every fixed
positive integer $k$.  In his indexing the coefficient called
$C_m(k)$ is our $C_{1,m+1}(k-1)$, which has degree exactly $m$.
The first numerical constant proposed there is identified exactly
in Corollary~\ref{cor:cohen-two-constants}.
The largest-root estimate, together with real-rootedness, also gives
Conjecture~7.1(1), and Proposition~\ref{prop:cohen-interlacing}
gives Conjecture~7.1(3).  These are statements about the precise
conjectures in the cited version; no exhaustive priority claim is
intended.  The derivative interpretation proposed for $k\leq0$
is not needed here.  The theorem is a fixed-$j$, fixed-$k$ result:
its remainder constants are not asserted uniform for growing indices.

The bulk law and the edge theorem describe different scales.
A positive fraction of the zeros stays in a bounded interval of the
logarithmic coordinate $X=L$, whereas every fixed number of largest
zeros grows linearly with $n$.  The $L^{-1}$ tail of the bulk limiting
measure agrees with the leading edge scale $n/j$.  A uniform
transition for $j\to\infty$ with $j/n\to0$, and a regime in which
$k$ also grows, remain natural research problems.
